\chapter{Technology Stack and Software Implementation}
\section{Frontend and backend}
React components implement the dashboard, file explorer, sharing, account, activity, cloud and Ask views. React Router handles navigation; Axios provides the client API layer; Tailwind utility classes and the existing design tokens support light/dark presentation. The dependency manifest does not establish React Query as the active state-management layer, despite its appearance in the proposal.

The Node.js backend uses Express middleware for authentication, validation, rate limits and error handling. Zod validates request shapes; bcrypt verifies password hashes; jsonwebtoken handles access tokens; cookie-parser supports refresh-cookie transport. Prisma migrations and generated clients define the relational interface. Vite builds the frontend; TypeScript compilation and linting are separate checks from runtime acceptance.
\par\begin{table}[htbp]\caption{Implemented technology stack and representative source locations}\small\begin{tabularx}{\linewidth}{p{.25\linewidth}Yp{.32\linewidth}}\toprule Technology / layer&Purpose&Source location\\\midrule
React, TypeScript, Vite&SPA interface and production build&\sourcepath{frontend/src}\\
Express, Zod, JWT, bcrypt&REST, validation and identity&\sourcepath{backend/src/modules}\\
PostgreSQL / Prisma&Ownership, folders, grants and jobs&\sourcepath{backend/prisma}\\
FastAPI / Pydantic&Authenticated private adapter&\sourcepath{intelligence/app/internal.py}\\
MiniLM / PyTorch / NumPy&384D normalised embeddings&\sourcepath{intelligence/app/embeddings.py}\\
SQLite / FTS5&Passages, vectors, BM25 and history&\sourcepath{intelligence/app/database.py}\\
Ollama / Qwen2.5 3B&Local structured generation&\sourcepath{intelligence/app/llm.py}\\
Docker Compose / Git&Runtime isolation and checkpoints&\sourcepath{deployment/}\\\bottomrule\end{tabularx}\end{table}\par
\section{Version and reproducibility policy}
The measured Linux stack uses Python 3.13.18 and Torch 2.15.0+cpu. The Node runtime observed in the workstation environment was 24.19.0; production images pin a Node 24 base digest. Package manifests contain version ranges, so a range such as React \texttt{\^{}18.2.0} is not presented as an installed-version measurement. The lockfiles and image digests are the reproducibility inputs. The 75-package Python lock and vendor CPU runtime archive are retained by the project.

The application does not implement Kubernetes, ECS, RDS, a load balancer or CloudFront in its accepted runtime. A separate original CDK design exists, but its presence is not evidence that those resources are active. The deployed serverless stack is reviewed through its actual CloudFormation template and worker source~\cite{cloudformation}.
\section{Implementation traceability}
Key source paths are recorded throughout the report and in Appendix~\ref{app:source}. Framework documentation explains the tool contracts~\cite{react,express,fastapi,prisma}; implementation claims are grounded in repository code and tests. This distinction also avoids attributing a framework's possible scalability to a measured property of this particular deployment.
\section{Generation-aware background work}
The application also retains a durable PostgreSQL indexing-job subsystem. Jobs carry a file generation, attempt count, run-after time and lease owner. Claiming uses row locking with \texttt{SKIP LOCKED}, and a completed older generation cannot replace current index state. Repeated crashes become an explicit failed job rather than an indefinitely running success indicator. This queue is distinct from the private engine's single active question-generation lock and from the AWS SQS ingestion queue.
Figure~\ref{fig:jobs} summarises the durable indexing-job lifecycle.
\diagram{job_lifecycle}{Durable application indexing-job lifecycle. Generation checks and bounded attempts prevent stale work from becoming the current result.}{jobs}
\FloatBarrier
